\section{Proposed Solution}

\subsection{Preliminary Architecture}

Figure~\ref{fig:context} shows the C4 System Context diagram of Room4U.
Students use Room4U to browse rooms, submit booking requests and track
their status. Library administrators use it to approve or reject requests,
cancel approved bookings and manage rooms and maintenance periods.
Firebase Authentication is the only external system: it verifies user
identities, while Room4U itself checks each user's role and permissions and
stores all business data. The database is an internal part of Room4U and
therefore does not appear as an external system at this level. Containers,
components and technical mechanisms are detailed in PA3.

\begin{figure}[H]
\centering
\hyphenpenalty=10000 \exhyphenpenalty=10000 \relax
\begin{tikzpicture}[
  font=\small,
  person/.style={draw=c4person, fill=c4person, text=white, rounded corners=5pt,
                 minimum width=4.2cm, minimum height=1.9cm,
                 text width=3.8cm, align=center, inner sep=4pt},
  system/.style={draw=c4system, fill=c4system, text=white, rounded corners=5pt,
                 minimum width=4.8cm, minimum height=2.2cm,
                 text width=4.4cm, align=center, inner sep=4pt},
  external/.style={draw=c4ext, fill=c4ext, text=white, rounded corners=5pt,
                   minimum width=4.4cm, minimum height=1.9cm,
                   text width=4.0cm, align=center, inner sep=4pt},
  rel/.style={-{Stealth[length=2.6mm]}, thick, shorten >=1pt, shorten <=1pt},
  lbl/.style={font=\footnotesize, text width=3.3cm, align=center,
              fill=white, inner sep=2pt}
]
\node[person]   (student) at (-4.6, 4.4) {\textbf{Student}\\{\footnotesize [Person]}\\{\footnotesize Looks for study rooms and requests bookings}};
\node[person]   (admin)   at ( 4.6, 4.4) {\textbf{Library Administrator}\\{\footnotesize [Person]}\\{\footnotesize Processes requests and manages rooms}};
\node[system]   (r4u)     at ( 0.0, 0.0) {\textbf{Room4U}\\{\footnotesize [Software System]}\\{\footnotesize Website for booking library study rooms}};
\node[external] (fb)      at ( 0.0,-4.0) {\textbf{Firebase Authentication}\\{\footnotesize [External System]}\\{\footnotesize Sign-in and identity service}};

\draw[rel] (student.south) -- node[lbl] {Searches rooms, submits and tracks booking requests} (r4u.north west);
\draw[rel] (admin.south)   -- node[lbl] {Approves or rejects requests, manages rooms and maintenance} (r4u.north east);
\draw[rel] (r4u.south)     -- node[lbl] {Verifies user identity tokens} (fb.north);
\end{tikzpicture}
\caption{C4 System Context diagram of Room4U. The database is an internal
part of Room4U and is therefore not drawn as an external system.}
\label{fig:context}
\end{figure}

\textbf{Preliminary business rules.} The rules below are working assumptions
used for scoping and estimation. They will be confirmed and refined in PA2.
\begin{itemize}[leftmargin=*,itemsep=2pt]
  \item The library opens from 08:00 to 20:00 every day. All times use the
        \texttt{Asia/Ho\_Chi\_Minh} time zone.
  \item A booking starts on a 30-minute boundary (a PA1 planning assumption
        to be confirmed and detailed in PA2), lasts between 30 and 120 minutes,
        and can be made at most 7 days in advance.
  \item A Pending request does not hold a room. When an administrator approves
        a request, the back end re-checks availability, and concurrent approvals
        are guarded by a transaction or locking mechanism designed in PA3.
  \item Two time intervals $[\mathit{start},\mathit{end})$ conflict when
        $\mathit{start}_1 < \mathit{end}_2$ and $\mathit{start}_2 < \mathit{end}_1$.
        A booking that ends at 10:00 therefore does not conflict with one that
        starts at 10:00.
  \item A student may cancel their own Pending request or Approved booking
        before its start time. Rejected and Cancelled bookings do not occupy
        the room.
  \item An administrator cannot create a maintenance period that overlaps an
        Approved booking; the affected bookings must first be cancelled with
        a reason. A suspended room does not accept new requests.
\end{itemize}

\subsection{Technology Stack}

\begin{table}[H]
\centering
\caption{Technology stack of Room4U.}
\label{tab:stack}
\begin{tabularx}{\textwidth}{@{}L{2.9cm}L{4.4cm}Y@{}}
\toprule
\textbf{Layer} & \textbf{Technology} & \textbf{Reason for choice}\\
\midrule
Front end & HTML, CSS, JavaScript &
  Suitable for a demo with a small number of screens; runs in any modern
  browser without extra tooling.\\
Back end & Node.js with Express &
  Same language as the front end; a layered structure (routes, services,
  repositories) keeps the business rules in one place.\\
Database & PostgreSQL &
  Rooms, users and bookings are relational data; transactions and
  constraints help prevent double booking.\\
Authentication & Firebase Authentication &
  Handles sign-in for pre-created accounts; the back end verifies tokens and
  checks the roles stored in the business database.\\
Hosting and deployment & Linux VPS with Nginx &
  Runs the Node.js back end and PostgreSQL; HTTPS, with configuration kept
  in environment variables.\\
Testing and CI & Vitest, Supertest, Playwright, GitHub Actions &
  Unit, API, and end-to-end tests are planned. Initial GitHub Actions CI setup
  is planned for W8, with checks added as implementation proceeds.\\
Source control and tasks & GitHub, GitHub Projects &
  GitHub Projects tracks WBS-linked tasks; the agreed workflow requires
  pull-request review before merge.\\
\bottomrule
\end{tabularx}
\end{table}

The stack is a planning decision rather than a final architecture choice.
During the design phase (Weeks 6--7) the team will validate it with a short
environment spike: running the front end, back end and database together on a
development machine, verifying a Firebase ID token in the Express back end, and
making a trial deployment to a Linux server. If any step fails, the
alternative will be recorded as an Architecture Decision Record in PA3.
Secrets such as API keys and service-account files will never be committed to
the repository.

\subsection{Hardware and Operating Environment}

Table~\ref{tab:env} lists the environments needed to develop, test and
demonstrate Room4U.

\begin{table}[H]
\centering
\caption{Development, test and demonstration environments.}
\label{tab:env}
\begin{tabularx}{\textwidth}{@{}L{3.3cm}Y@{}}
\toprule
\textbf{Environment} & \textbf{Planning assumption}\\
\midrule
Development machines &
  Laptops running Windows, macOS or Linux with at least 8~GB of RAM, Git,
  Visual Studio Code and one agreed Node.js LTS version.\\
Demo server &
  Linux server with 2 vCPUs, 4~GB of RAM and 20~GB of storage, hosting one
  back end instance and one database. These figures are assumptions and will
  be checked with a smoke test.\\
User devices &
  Laptops or phones with a modern browser. Testing focuses on Chrome and Edge
  on desktop, plus one mobile viewport.\\
Test data &
  Eight simulated rooms; pre-created Student and Administrator accounts; data
  sets covering schedule conflicts, time limits, capacity limits and
  permissions.\\
Demo operation &
  HTTPS, Internet access, a Firebase test project, a data backup before the
  demo, and a backup video in case the network fails.\\
\bottomrule
\end{tabularx}
\end{table}